\documentclass[11pt,a4paper,oneside,fontset=windows]{ctexbook}
\usepackage{amsthm,amssymb} % math packages
\usepackage{amsmath}
\usepackage{mathrsfs}
\usepackage{bm} % bold math
\usepackage[hidelinks]{hyperref} % hide ref in table of contents
\usepackage{pifont} % \ding{73} makes a beautiful star
\usepackage{tikz}
\usepackage{fontspec}
\usepackage{color}
\usepackage{geometry}
\usepackage{stmaryrd}
\usepackage{subfiles}
\usepackage{pgfplots}
\usepackage{pgfornament}
\usepackage{enumerate}
\usepackage{mdwlist}
\usepackage{booktabs}
\usepackage{ulem}
\usepackage[level]{datetime}


\allowdisplaybreaks[4]

\setcounter{tocdepth}{1}
\definecolor{navy}{rgb}{0,0,0.5}
% \tag{} numbers the equation
\def\d#1{\text{d}#1}
\def\D#1{\text{d}(#1)}
\def\p#1{\partial#1}
\def\ds#1#2{\frac{\text{d}#1}{\text{d}#2}}
\def\pds#1#2{\frac{\partial#1}{\partial#2}}
\def\jishu#1#2{\sum_{#1=1}^{+\infty}#2}
\def\rlinsp#1{\mathbb{R}^{#1}}
\def\clinsp#1{\mathbb{C}^{#1}}
\def\ntoinf#1{\lim\limits_{#1\rightarrow\infty}}
\def\tinysec#1{\noindent\textbf{#1}\indent}
\def\concept#1{\textbf{\uline{#1}}}
\geometry{a4paper, scale=0.75}

\newenvironment{solution}{\noindent{\kaishu 解}.}{  }

\title{\Huge{\heiti 数学支配着宇宙}\\Math Dominates the Universe}
\author{Jingze Xing}

\pgfplotsset{compat=1.16} % advised setting
\begin{document}
%\ThisCenterWallPaper{}{back.jpg}
\thispagestyle{empty}
\begin{tikzpicture}[color=navy,overlay,remember picture,
    every node/.style={inner sep=5pt}]
    \node[anchor=north west] at (current page.north west)
    {\pgfornament[width=3cm]{63}};
    \node[anchor=north east] at (current page.north east)
    {\pgfornament[width=3cm,symmetry=v]{63}};
    \node[anchor=south west] at (current page.south west)
    {\pgfornament[width=3cm,symmetry=h]{63}};
    \node[anchor=south east] at (current page.south east)
    {\pgfornament[width=3cm,symmetry=c]{63}};
\end{tikzpicture}
\begin{figure}
    \centering
    \begin{minipage}[t][10cm][b]{\textwidth}
        \centering
        \begin{tikzpicture}[color=black,remember picture]
            \node[minimum width=12cm, minimum height=3.6cm](vecbox){};
            \node[anchor=north] at (vecbox.north)
            {\pgfornament[height=0.27cm]{89}};
            \node[anchor=south] at (vecbox.south)
            {\pgfornament[height=0.27cm]{89}};
            
            \node[inner sep=10pt,align=center] (text) at (vecbox.center){\Huge\textbf{Math Dominates the Universe}\\[1em]\Huge{\heiti 数学支配着宇宙}};
        \end{tikzpicture}
    \end{minipage}
    \centering
    \begin{minipage}[t][3.5cm][b]{\textwidth}
        \centering
        \begin{tikzpicture}[color=black,remember picture]
            \node[minimum width=8.6cm, minimum height=2cm](vecbox){};
            \node[inner sep=0pt,align=center,anchor=north] (text) at ([yshift={-7pt}]vecbox.north){\LARGE\textbf{Jingze Xing}};
            \node[inner sep=0pt,align=center,anchor=south] (text) at (vecbox.south){\Large{\today}};
            \node[inner sep=0pt,anchor=north west] at (vecbox.north west)
            {\pgfornament[height=1cm]{113}};
            \node[inner sep=0pt,anchor=north east] at (vecbox.north east)
            {\pgfornament[height=1cm,symmetry=v]{113}};
        \end{tikzpicture}
    \end{minipage}
\end{figure}
\tableofcontents
\newpage

\subfile{1-calculus1.tex}
\subfile{2-calculus2.tex}
\subfile{7-linear-algebra.tex}
% \subfile{3-discrete.tex}
\subfile{4-complex-functions.tex}
\subfile{5-probability.tex}
\subfile{6-stochastic-process.tex}

\subsection{特征值与特征向量}
\subsubsection{练习5.2.23}
给定$m$阶方阵$A_1$，$n$阶方阵$A_2$和$m\times n$矩阵$B$，证明若$A_1$和$A_2$没有相同的特征值，则关于$m\times n$矩阵$X$的矩阵方程$A_1X-XA_2=B$有唯一解.\\
\begin{proof}
  因为$X$与$B$规格相同，所以只要证明该变换是单射即可证明是双射.\\
  即证：$A_1X=XA_2$仅有零解
  \begin{align*}
    A_1^2X=(A_1X) & A_2=XA_2^2                      \\
    A_1^3X=(A_1X) & A_2^2=XA_2^3                    \\
                  & \vdots                          \\
    A_1^nX=(A_1X) & A_2^{n-1}=XA_2^n\tag{\ding{72}}
  \end{align*}
  考虑上述等式的线性组合，则
  $$f(A_1)X=Xf(A_2)$$
  恒成立.\\
  不妨令$f=p_A$，则$f(A_1)=O$.令$\lambda_i$为$A_1$的特征值，则$\det(\lambda_iI_n-A_2)\neq 0$，所以
  \begin{align*}
    \det(f(A_2))=\prod\limits_{i=1}^{n}{(\lambda_i I_n-A_2)}\neq 0 \tag{\ding{72}}
  \end{align*}
  故$f(A_2)$可逆，因此$$X=f(A_1)X\left[f(A_2)\right]^{-1}=O$$
\end{proof}
\subsection{实对称矩阵}
对实对称矩阵，$A$正定\\$\iff$$A$的特征值都是正数\\$\iff$$A=TT^{\mathrm{T}}$,$T$可逆\\$\iff$$A=LDL^{\mathrm{T}}$\\$\iff$$A$的顺序主子式都是正数

\end{document}